\section{Formulario}

\subsection{Flujo eléctrico}

Para una superficie plana de área \(A\) en un campo eléctrico uniforme,

\[
\boxed{\Phi_E=EA\cos\theta}
\]

donde \(\theta\) es el ángulo entre el campo eléctrico y la normal a la superficie.

Casos particulares:

\[
\theta=0^\circ
\quad\Longrightarrow\quad
\boxed{\Phi_E=EA}
\]

\[
\theta=90^\circ
\quad\Longrightarrow\quad
\boxed{\Phi_E=0}
\]

\[
\theta=180^\circ
\quad\Longrightarrow\quad
\boxed{\Phi_E=-EA}.
\]

Unidad SI:

\[
\boxed{
[\Phi_E]=\si{\newton\metre\squared\per\coulomb}
}.
\]


\subsection{Signo del flujo}

Para una superficie cerrada,

\[
\boxed{\text{campo saliente}\Longrightarrow\Phi_E>0}
\]

\[
\boxed{\text{campo entrante}\Longrightarrow\Phi_E<0}.
\]

Un flujo nulo no implica necesariamente un campo nulo:

\[
\boxed{\Phi_E=0\;\not\Rightarrow\;E=0}.
\]


\subsection{Flujo sobre una superficie general}

Para un elemento de área \(dA\),

\[
\boxed{d\Phi_E=E\cos\theta\,dA}.
\]

Para una superficie abierta,

\[
\boxed{
\Phi_E=\int_S\vec E\cdot d\vec A
}
\]

o, equivalentemente,

\[
\boxed{
\Phi_E=\int_S E\cos\theta\,dA
}.
\]


\subsection{Flujo sobre una superficie cerrada}

\[
\boxed{
\Phi_E=\oint_S\vec E\cdot d\vec A
}
\]

En una superficie cerrada, \(d\vec A\) se define siempre hacia el exterior.


\subsection{Ley de Gauss}

\[
\boxed{
\oint_S\vec E\cdot d\vec A
=
\frac{Q_{\text{enc}}}{\varepsilon_0}
}
\]

Por tanto,

\[
\boxed{
\Phi_E=\frac{Q_{\text{enc}}}{\varepsilon_0}
}.
\]

Permitividad eléctrica del vacío:

\[
\boxed{
\varepsilon_0\approx
\SI{8.854e-12}{\coulomb\squared\per\newton\metre\squared}
}.
\]

Relación con la constante de Coulomb:

\[
\boxed{
k=\frac{1}{4\pi\varepsilon_0}
}.
\]


\subsection{Carga neta encerrada}

Para cargas puntuales,

\[
\boxed{
Q_{\text{enc}}=\sum_iq_i
}.
\]

La suma es algebraica y conserva el signo de cada carga.

Si

\[
Q_{\text{enc}}>0,
\]

entonces

\[
\boxed{\Phi_E>0}.
\]

Si

\[
Q_{\text{enc}}<0,
\]

entonces

\[
\boxed{\Phi_E<0}.
\]

Si

\[
Q_{\text{enc}}=0,
\]

entonces

\[
\boxed{\Phi_E=0}.
\]


\subsection{Cargas externas}

Las cargas situadas fuera de una superficie cerrada no forman parte de
\(Q_{\text{enc}}\).

Su contribución al flujo neto es

\[
\boxed{
\Phi_{E,\text{externa}}=0
}.
\]

Esto no significa que su campo sobre la superficie sea nulo.


\subsection{Distribuciones continuas y carga encerrada}

Para una distribución lineal,

\[
\boxed{
Q_{\text{enc}}=\int\lambda\,d\ell
}.
\]

Si \(\lambda\) es constante,

\[
\boxed{Q_{\text{enc}}=\lambda L}.
\]

Para una distribución superficial,

\[
\boxed{
Q_{\text{enc}}=\int\sigma\,dA
}.
\]

Si \(\sigma\) es constante,

\[
\boxed{Q_{\text{enc}}=\sigma A}.
\]

Para una distribución volumétrica,

\[
\boxed{
Q_{\text{enc}}=\int\rho\,dV
}.
\]

Si \(\rho\) es constante,

\[
\boxed{Q_{\text{enc}}=\rho V}.
\]


\subsection{Simetría esférica}

Para una superficie gaussiana esférica de radio \(r\),

\[
\boxed{A=4\pi r^2}.
\]

Si el campo tiene magnitud constante sobre la superficie y es radial,

\[
\boxed{\Phi_E=E(4\pi r^2)}.
\]

La ley de Gauss queda

\[
\boxed{
E(4\pi r^2)
=
\frac{Q_{\text{enc}}}{\varepsilon_0}
}.
\]

Por tanto,

\[
\boxed{
E=
\frac{|Q_{\text{enc}}|}
{4\pi\varepsilon_0r^2}
}
\]

cuando la simetría permite tratar \(Q_{\text{enc}}\) de esta forma.


\subsection{Carga puntual}

Para una carga puntual \(Q\),

\[
\boxed{
E=k\frac{|Q|}{r^2}
}
\]

con

\[
\boxed{
k=\frac{1}{4\pi\varepsilon_0}
}.
\]

Si

\[
Q>0,
\]

el campo apunta radialmente hacia afuera.

Si

\[
Q<0,
\]

el campo apunta radialmente hacia la carga.


\subsection{Cáscara esférica uniformemente cargada}

Para una cáscara esférica delgada de radio \(R\) y carga total \(Q\):

Exterior,

\[
\boxed{
E=k\frac{|Q|}{r^2}
\qquad (r>R)
}.
\]

Interior,

\[
\boxed{
E=0
\qquad (r<R)
}.
\]

En la superficie,

\[
\boxed{
E(R)=k\frac{|Q|}{R^2}
}.
\]


\subsection{Esfera sólida uniformemente cargada}

Para una esfera no conductora de radio \(R\) y carga total \(Q\),

\[
\boxed{
\rho=\frac{Q}{\frac43\pi R^3}
}.
\]

Para

\[
r<R,
\]

la carga encerrada es

\[
\boxed{
Q_{\text{enc}}
=
Q\frac{r^3}{R^3}
}.
\]

El campo interior es

\[
\boxed{
E=k\frac{|Q|}{R^3}r
\qquad (r<R)
}.
\]

Por tanto,

\[
\boxed{E\propto r\qquad(r<R)}.
\]

Para

\[
r\geq R,
\]

\[
\boxed{
E=k\frac{|Q|}{r^2}
}.
\]

En el centro,

\[
\boxed{E(0)=0}.
\]

En la superficie,

\[
\boxed{
E(R)=k\frac{|Q|}{R^2}
}.
\]


\subsection{Simetría cilíndrica}

Para una superficie gaussiana cilíndrica de radio \(r\) y longitud \(L\),

\[
\boxed{
A_{\text{lateral}}=2\pi rL
}.
\]

Si el campo es radial,

\[
\boxed{
\Phi_{\text{lateral}}=E(2\pi rL)
}.
\]

Si el campo es paralelo a las tapas,

\[
\boxed{
\Phi_{\text{tapas}}=0
}.
\]


\subsection{Línea infinita uniformemente cargada}

Para una línea infinita con densidad lineal \(\lambda\),

\[
\boxed{
Q_{\text{enc}}=\lambda L
}.
\]

Aplicando Gauss,

\[
E(2\pi rL)
=
\frac{\lambda L}{\varepsilon_0}.
\]

Por tanto,

\[
\boxed{
E=\frac{|\lambda|}{2\pi\varepsilon_0r}
}
\]

o

\[
\boxed{
E=\frac{2k|\lambda|}{r}
}.
\]

Dependencia con la distancia:

\[
\boxed{
E\propto\frac{1}{r}
}.
\]

Si

\[
\lambda>0,
\]

el campo apunta alejándose de la línea.

Si

\[
\lambda<0,
\]

el campo apunta hacia la línea.


\subsection{Simetría plana}

Para una lámina infinita uniformemente cargada, se utiliza una superficie
gaussiana cilíndrica que atraviesa la lámina.

El flujo por la superficie lateral es

\[
\boxed{\Phi_{\text{lateral}}=0}.
\]

Si cada tapa tiene área \(A\),

\[
\boxed{\Phi_E=2EA}.
\]


\subsection{Lámina infinita uniformemente cargada}

Para una densidad superficial \(\sigma\),

\[
\boxed{
Q_{\text{enc}}=\sigma A
}.
\]

Aplicando la ley de Gauss,

\[
2EA=\frac{\sigma A}{\varepsilon_0}.
\]

Por tanto,

\[
\boxed{
E=\frac{|\sigma|}{2\varepsilon_0}
}.
\]

El campo es independiente de la distancia:

\[
\boxed{E=\text{constante respecto de la distancia}}.
\]

Si

\[
\sigma>0,
\]

el campo apunta alejándose de la lámina.

Si

\[
\sigma<0,
\]

el campo apunta hacia ella.


\subsection{Comparación de dependencias con la distancia}

Carga puntual o distribución esférica exterior:

\[
\boxed{
E\propto\frac{1}{r^2}
}.
\]

Línea infinita:

\[
\boxed{
E\propto\frac{1}{r}
}.
\]

Lámina infinita:

\[
\boxed{
E\propto r^0
}
\]

es decir,

\[
\boxed{E=\text{constante}}.
\]


\subsection{Condiciones útiles para simplificar el flujo}

Si el campo es perpendicular a la superficie,

\[
\theta=0^\circ
\quad\Longrightarrow\quad
\boxed{d\Phi_E=E\,dA}.
\]

Si el campo es paralelo a la superficie,

\[
\theta=90^\circ
\quad\Longrightarrow\quad
\boxed{d\Phi_E=0}.
\]

Si \(E\) es constante sobre una región de área \(A\),

\[
\boxed{
\int E\,dA=EA
}.
\]


\subsection{Superficies gaussianas frecuentes}

\[
\boxed{
\begin{array}{c|c}
\text{Simetría} & \text{Superficie gaussiana conveniente}\\
\hline
\text{Esférica} & \text{Esfera concéntrica}\\
\text{Cilíndrica} & \text{Cilindro coaxial}\\
\text{Plana} & \text{Cilindro corto o pastilla}
\end{array}
}
\]


\subsection{Procedimiento para calcular flujo}

\begin{enumerate}
    \item Identificar la superficie.
    \item Determinar el campo \(E\).
    \item Determinar el área \(A\).
    \item Identificar el ángulo entre el campo y la normal.
    \item Aplicar
    \[
    \Phi_E=EA\cos\theta
    \]
    si \(E\) y \(\theta\) son constantes.
    \item Si cambian sobre la superficie, utilizar
    \[
    \Phi_E=\int E\cos\theta\,dA.
    \]
    \item Asignar correctamente el signo del flujo.
\end{enumerate}


\subsection{Procedimiento para aplicar la ley de Gauss}

\begin{enumerate}
    \item Identificar la simetría de la distribución.
    \item Determinar cualitativamente la dirección del campo.
    \item Elegir una superficie gaussiana cerrada compatible con la simetría.
    \item Identificar qué partes de la superficie aportan flujo.
    \item Calcular \(Q_{\text{enc}}\).
    \item Aplicar
    \[
    \oint_S\vec E\cdot d\vec A
    =
    \frac{Q_{\text{enc}}}{\varepsilon_0}.
    \]
    \item Simplificar la integral mediante la simetría.
    \item Despejar \(E\).
    \item Determinar el sentido del campo según el signo de la carga.
\end{enumerate}


\subsection{Advertencias}

\begin{itemize}
    \item En
    \[
    \Phi_E=EA\cos\theta,
    \]
    \(\theta\) se mide respecto de la normal, no respecto del plano.

    \item La ley de Gauss requiere una superficie cerrada.

    \item Sólo las cargas encerradas forman parte de \(Q_{\text{enc}}\).

    \item Las cargas externas pueden producir campo sobre la superficie, aunque su
    contribución al flujo neto sea cero.

    \item
    \[
    \Phi_E=0
    \]
    no implica
    \[
    E=0.
    \]

    \item La ley de Gauss siempre es válida, pero sólo permite calcular \(E\)
    fácilmente cuando existe suficiente simetría.

    \item Una superficie gaussiana es imaginaria; no tiene que coincidir con una
    superficie física.
\end{itemize}


\subsection{Resumen de fórmulas}

\begin{align*}
\Phi_E &= EA\cos\theta,\\
\Phi_E &= \int_S\vec E\cdot d\vec A,\\
\Phi_E &= \oint_S\vec E\cdot d\vec A,\\
\Phi_E &= \frac{Q_{\text{enc}}}{\varepsilon_0},\\
Q_{\text{enc}} &= \sum_iq_i,\\
Q_{\text{enc}} &= \int\lambda\,d\ell,\\
Q_{\text{enc}} &= \int\sigma\,dA,\\
Q_{\text{enc}} &= \int\rho\,dV,\\
E_{\text{esférica}}(4\pi r^2) &= \frac{Q_{\text{enc}}}{\varepsilon_0},\\
E_{\text{puntual}} &= k\frac{|Q|}{r^2},\\
E_{\text{cáscara,int}} &= 0,\\
E_{\text{esfera,int}} &= k\frac{|Q|}{R^3}r,\\
E_{\text{línea}} &= \frac{|\lambda|}{2\pi\varepsilon_0r},\\
E_{\text{lámina}} &= \frac{|\sigma|}{2\varepsilon_0}.
\end{align*}